% META: {"id": "1SPE-SUITES-COURS-01", "chapitre": "1SPE-SUITES", "type_objet": "cours", "sous_type": "diagnostic", "titre": "Diagnostic d'entree", "temps_gabarit": 2, "duree_min": 20, "prerequis_testes": ["R1", "R2", "R3", "R4", "R5"], "mode_creation": "ex_nihilo", "sources_inspiration": [], "status": "generated"}
% BEGIN-VERIFY
% from sympy import symbols, factor, expand, Rational, solve, Symbol
% x = symbols('x')
% # Q1 : développer (3x - 2)(x + 5)
% expr_q1 = expand((3*x - 2)*(x + 5))
% assert expr_q1 == 3*x**2 + 13*x - 10
% # Q2 : factoriser 4x^2 - 9
% expr_q2 = factor(4*x**2 - 9)
% assert expr_q2 == (2*x - 3)*(2*x + 3)
% # Q3 : coefficient multiplicateur d'une baisse de 15 %
% cm_q3 = Rational(85, 100)
% assert cm_q3 == Rational(17, 20)
% # Q4 : capital apres deux evolutions successives de +5 % puis -5 %
% capital_depart = 2000
% capital_final = capital_depart * Rational(105, 100) * Rational(95, 100)
% assert capital_final == Rational(199500, 100)   # 1995
% # Q5 : image de x=3 par f(x) = 2x^2 - 5
% def f(val): return 2*val**2 - 5
% assert f(3) == 13
% # Q6 : antecedent de 7 par g(x) = 3x + 1
% x_var = Symbol('x')
% sol_q6 = solve(3*x_var + 1 - 7, x_var)
% assert sol_q6 == [2]
% # Q7 : signe de (x - 1)(x + 3) pour x = 4
% val_q7 = (4 - 1)*(4 + 3)
% assert val_q7 > 0
% # Q8 : signe de -2x + 6 pour x = 5
% val_q8 = -2*5 + 6
% assert val_q8 < 0
% END-VERIFY

\section*{Temps 2 — Diagnostic d'entrée \hfill \normalfont\small(15–20 min, sans calculatrice)}

Réponds à chaque question sur ta copie. Ne cherche pas à aller vite : l'objectif est de repérer ce que tu maîtrises, pas de tout réussir.

\bigskip

\subsection*{R1 — Calcul littéral}

\begin{enumerate}

  \item Développe et réduis $(3x - 2)(x + 5)$.

  \item Factorise $4x^2 - 9$.

\end{enumerate}

\subsection*{R2 — Pourcentages}

\begin{enumerate}
  \setcounter{enumi}{2}

  \item Le prix d'un article baisse de $15$~\%. Quel est le coefficient multiplicateur correspondant à cette évolution ?

  \item Un capital de \nombre{2000}~€ subit deux évolutions successives : $+5$~\% puis $-5$~\%. Quel est le capital final ? Est-ce le même qu'après une évolution globale de $0$~\% ? Justifie.

\end{enumerate}

\subsection*{R3 — Fonctions}

\begin{enumerate}
  \setcounter{enumi}{4}

  \item Soit $f$ la fonction définie sur $\mathbb{R}$ par $f(x) = 2x^2 - 5$. Calcule $f(3)$.

  \item Soit $g$ la fonction définie sur $\mathbb{R}$ par $g(x) = 3x + 1$. Détermine l'antécédent de $7$ par $g$.

\end{enumerate}

\subsection*{R4 — Inégalités}

\begin{enumerate}
  \setcounter{enumi}{6}

  \item Étudie le signe de $(x - 1)(x + 3)$ pour $x = 4$, puis pour $x = -2$.

  \item Pour quelles valeurs de $x$ a-t-on $-2x + 6 < 0$ ?

\end{enumerate}

\subsection*{R5 — Python}

\begin{enumerate}
  \setcounter{enumi}{8}

  \item On exécute le programme suivant :

\begin{python}
s = 0
for k in range(1, 5):
    s = s + k
print(s)
\end{python}

  Quelle valeur est affichée ? Explique le rôle de la variable \texttt{s}.

  \item On exécute le programme suivant :

\begin{python}
u = 3
for i in range(4):
    u = 2 * u + 1
print(u)
\end{python}

  Complète le tableau de valeurs en suivant l'exécution pas à pas, puis indique la valeur affichée.

  \begin{center}
  \begin{tabular}{|c|c|}
  \hline
  Valeur de \texttt{i} & Valeur de \texttt{u} après la ligne \texttt{u = 2*u+1} \\
  \hline
  $0$ & \phantom{xxx} \\
  \hline
  $1$ & \phantom{xxx} \\
  \hline
  $2$ & \phantom{xxx} \\
  \hline
  $3$ & \phantom{xxx} \\
  \hline
  \end{tabular}
  \end{center}

\end{enumerate}

\bigskip

\subsection*{Grille de décodage}

Reporte tes résultats dans ce tableau, puis suis les conseils de remédiation avant de commencer le cours.

\begin{center}
\begin{tabular}{|c|c|l|l|}
\hline
\textbf{Questions} & \textbf{Prérequis testé} & \textbf{Si tu as échoué...} & \textbf{Fiche à lire} \\
\hline
1 et 2 & R1 — Calcul littéral & ... avant de commencer C1 & Fiche R1 \\
\hline
3 et 4 & R2 — Pourcentages & ... avant de commencer C3, C6 & Fiche R2 \\
\hline
5 et 6 & R3 — Fonctions & ... avant de commencer C1, C5 & Fiche R3 \\
\hline
7 et 8 & R4 — Inégalités & ... avant de commencer C5 & Fiche R4 \\
\hline
9 et 10 & R5 — Python & ... avant de commencer C7 & Fiche R5 \\
\hline
\end{tabular}
\end{center}

\bigskip

Si tu as réussi toutes les questions sans difficulté, tu peux démarrer directement le Temps~3.
